\documentclass[11pt]{article}
\usepackage[margin=1in]{geometry}
\usepackage{amsmath,amssymb,amsthm}
\usepackage{mathrsfs}
\usepackage{graphicx}
\usepackage[hidelinks]{hyperref}

\newtheorem{theorem}{Theorem}
\newtheorem{lemma}{Lemma}
\newtheorem{corollary}{Corollary}
\newtheorem{remark}{Remark}

\newcommand{\E}{\mathbb{E}}
\newcommand{\R}{\mathbb{R}}
\newcommand{\D}{\mathbb{D}}
\newcommand{\g}{\gamma}
\newcommand{\cL}{\mathcal L}
\newcommand{\1}{\mathbf 1}

\begin{document}

\begin{center}
{\Large A basis-free Lipschitz Malliavin chain rule}\\[4pt]
{\normalsize Full result for arXiv:1204.2946, Remark 3.11}
\end{center}

\section*{Source Question}

Pronk and Veraar prove in \cite[Proposition 3.10]{PV} that if \(E_0\) has a
Schauder basis, \(E_1\) is UMD, \(p\in(1,\infty)\), and
\(\phi:E_0\to E_1\) is Lipschitz, then
\[
  F\in \D^{1,p}(E_0) \quad\Longrightarrow\quad
  \phi(F)\in \D^{1,p}(E_1),
\]
and there is a bounded linear operator
\[
  T_F:\g(H,E_0)\to L^\infty(\Omega;\g(H,E_1))
\]
such that \(D(\phi(F))=T_F(DF)\).  In Remark 3.11 they write that they do not
know whether the assumption that \(E_0\) has a basis can be avoided.

\begin{center}
\includegraphics[width=0.95\linewidth]{figures/open_problem_crop.png}
\end{center}

\section*{Result}

If \(T:G\to L^\infty(\Omega;Y)\) is bounded and \(1\leq p<\infty\), write
\(T^{[p]}:L^p(\Omega;G)\to L^p(\Omega;Y)\) for the canonical amplification
defined on simple random variables by
\[
  T^{[p]}\Big(\sum_{j=1}^m \1_{A_j}R_j\Big)
  =\sum_{j=1}^m \1_{A_j}TR_j
\]
and then extended by continuity.  This is the meaning of \(T_F(DF)\) used
below, and is the one employed in the proof of \cite[Proposition 3.10]{PV}.

\begin{theorem}\label{thm:main}
Let \(H\) be a separable Hilbert space, let \(p\in(1,\infty)\), let \(E_0\) be
an arbitrary Banach space, and let \(E_1\) be a UMD Banach space.  If
\(\phi:E_0\to E_1\) is Lipschitz with constant \(L\), then for every
\(F\in\D^{1,p}(E_0)\) one has
\[
  \phi(F)\in \D^{1,p}(E_1).
\]
Moreover
\[
  \|D\phi(F)(\omega)\|_{\g(H,E_1)}
  \leq L\|DF(\omega)\|_{\g(H,E_0)}
  \quad\hbox{for a.e. }\omega\in\Omega ,
\]
and there exists a bounded linear operator
\[
  T_F:\g(H,E_0)\to L^\infty(\Omega;\g(H,E_1))
\]
with \(\|T_F\|\leq L\) such that
\[
  D(\phi(F))=T_F^{[p]}(DF).
\]
\end{theorem}

\begin{corollary}
The Schauder-basis assumption in \cite[Proposition 3.10]{PV} can be removed.
No separability or approximation property of \(E_0\) is needed.
\end{corollary}

\section*{Proof Intuition}

The source proof uses a basis to approximate the identity on \(E_0\) by
uniformly bounded finite-dimensional projections.  This is stronger than the
chain rule needs.  A Malliavin-Sobolev random variable is already obtained as a
limit of smooth random variables with finite-dimensional ranges.  We can apply
the finite-dimensional Lipschitz chain rule to those approximants and pass to
the limit.

The remaining operator \(T_F\) is only a factorization through the particular
random derivative \(DF\).  Once the pointwise estimate
\(\|D\phi(F)\|\leq L\|DF\|\) is known, the strong measurability of \(DF\) puts
its essential range in a separable subspace of \(\g(H,E_0)\).  On that
separable subspace we build a measurable scalar norming map.  The standard
Goodner--Kelley--Nachbin injectivity theorem for scalar \(L^\infty\) spaces
extends this scalar map to all of \(\g(H,E_0)\).  Multiplication by the
normalized derivative \(D\phi(F)/\|DF\|\) gives the desired \(T_F\).

\section*{A finite-dimensional lemma}

\begin{lemma}\label{lem:finite}
Let \(V\) be a finite-dimensional Banach space, let \(E_1\) be UMD, and let
\(\psi:V\to E_1\) be Lipschitz with constant \(L\).  If
\(G\in\D^{1,p}(V)\), then \(\psi(G)\in\D^{1,p}(E_1)\) and
\[
  \|D\psi(G)(\omega)\|_{\g(H,E_1)}
  \leq L\|DG(\omega)\|_{\g(H,V)}
  \quad\hbox{a.s.}
\]
\end{lemma}

\begin{proof}
Identify \(V\) linearly with \(\R^m\) and convolve \(\psi\) with a standard
smooth mollifier on \(V\).  This gives \(C^\infty\) maps
\(\psi_\varepsilon:V\to E_1\) such that
\[
  \|\psi_\varepsilon(x)-\psi(x)\|\to0\quad (x\in V),
  \qquad
  \|\psi_\varepsilon'(x)\|_{\cL(V,E_1)}\leq L .
\]
The smooth chain rule \cite[Proposition 3.9]{PV} gives
\[
  \psi_\varepsilon(G)\in\D^{1,p}(E_1),\qquad
  D\psi_\varepsilon(G)=\psi_\varepsilon'(G)DG,
\]
and the ideal property of \(\gamma\)-radonifying operators gives the pointwise
bound
\[
  \|D\psi_\varepsilon(G)(\omega)\|_{\g(H,E_1)}
  \leq L\|DG(\omega)\|_{\g(H,V)} .
\]
Also \(\psi_\varepsilon(G)\to\psi(G)\) in \(L^p(\Omega;E_1)\).  Since \(E_1\)
is UMD, \(\g(H,E_1)\) is reflexive, so the compactness lemma
\cite[Lemma 3.8]{PV} implies \(\psi(G)\in\D^{1,p}(E_1)\), after passing to a
weakly convergent derivative subsequence.

Mazur's lemma gives convex combinations of the derivatives
\(D\psi_{\varepsilon_j}(G)\) which converge strongly in
\(L^p(\Omega;\g(H,E_1))\) to \(D\psi(G)\).  The same pointwise bound is
preserved by convex combinations.  Passing to an almost surely convergent
subsequence yields the asserted pointwise estimate.
\end{proof}

\section*{The membership and pointwise bound}

Choose smooth random variables \(F_n\in\mathscr S\otimes E_0\) such that
\[
  F_n\to F \quad\hbox{in }\D^{1,p}(E_0).
\]
Each \(F_n\) takes values in a finite-dimensional subspace \(V_n\subseteq E_0\).
Applying Lemma \ref{lem:finite} to \(\phi|_{V_n}:V_n\to E_1\) gives
\[
  \phi(F_n)\in\D^{1,p}(E_1),\qquad
  \|D\phi(F_n)(\omega)\|_{\g(H,E_1)}
  \leq L\|DF_n(\omega)\|_{\g(H,E_0)} .
\]
Since \(\phi\) is Lipschitz,
\[
  \phi(F_n)\to\phi(F) \quad\hbox{in } L^p(\Omega;E_1).
\]
The derivatives \(D\phi(F_n)\) are bounded in
\[
  L^p(\Omega;\g(H,E_1)),
\]
so the compactness lemma again gives
\(\phi(F)\in\D^{1,p}(E_1)\) and, after passing to a subsequence,
\[
  D\phi(F_n)\rightharpoonup D\phi(F)
  \quad\hbox{in }L^p(\Omega;\g(H,E_1)).
\]

It remains to keep the pointwise estimate.  By Mazur's lemma there are convex
combinations
\[
  Z_k=\sum_{n\geq k} a_{k,n}D\phi(F_n)
\]
which converge strongly in \(L^p(\Omega;\g(H,E_1))\) to \(D\phi(F)\).  Set
\[
  M_k=\sum_{n\geq k} a_{k,n}\|DF_n\|_{\g(H,E_0)} .
\]
Since \(DF_n\to DF\) in \(L^p(\Omega;\g(H,E_0))\), we have
\[
  M_k\to \|DF\|_{\g(H,E_0)}
  \quad\hbox{in }L^p(\Omega).
\]
For every \(k\),
\[
  \|Z_k(\omega)\|_{\g(H,E_1)}\leq L M_k(\omega)
  \quad\hbox{a.s.}
\]
Taking an almost surely convergent subsequence of \(Z_k\) and \(M_k\) gives
\[
  \|D\phi(F)(\omega)\|_{\g(H,E_1)}
  \leq L\|DF(\omega)\|_{\g(H,E_0)}
\]
for almost every \(\omega\).

\section*{A factorization lemma}

\begin{lemma}\label{lem:factor}
Let \(G\) and \(Y\) be Banach spaces, let \(1\leq p<\infty\), and let
\(X\in L^p(\Omega;G)\), \(Z\in L^p(\Omega;Y)\) be strongly measurable.  Assume
\[
  \|Z(\omega)\|_Y\leq L\|X(\omega)\|_G
  \quad\hbox{for a.e. }\omega .
\]
Then there is a bounded linear operator
\[
  T:G\to L^\infty(\Omega;Y)
\]
with \(\|T\|\leq L\) such that \(T^{[p]}X=Z\).
\end{lemma}

\begin{proof}
By strong measurability there is a closed separable subspace \(S\subseteq G\)
such that \(X(\omega)\in S\) for almost every \(\omega\).  On \(S\) choose a
Borel norming selection
\[
  J:S\to B_{S^*},\qquad Jx(x)=\|x\|_S .
\]
This is the standard Kuratowski--Ryll-Nardzewski selection applied to
\(x\mapsto\{x^*\in B_{S^*}:x^*(x)=\|x\|\}\); separability makes the weak-star
compact ball \(B_{S^*}\) metrizable on the relevant coordinates.

Define \(A:S\to L^\infty(\Omega)\) by
\[
  (AR)(\omega)=J(X(\omega))(R),\qquad R\in S.
\]
Then \(A\) is linear and \(\|A\|\leq1\).  The scalar space
\(L^\infty(\Omega)\) is \(1\)-injective by the Goodner--Kelley--Nachbin
extension theorem, so \(A\) has an extension
\[
  \widetilde A:G\to L^\infty(\Omega)
\]
with \(\|\widetilde A\|\leq1\).

Put
\[
  W(\omega)=
  \begin{cases}
    Z(\omega)/\|X(\omega)\|_G, & X(\omega)\neq0,\\[3pt]
    0, & X(\omega)=0 .
  \end{cases}
\]
Then \(W\in L^\infty(\Omega;Y)\) and \(\|W\|_\infty\leq L\).  Multiplication by
\(W\) defines a bounded map
\[
  M_W:L^\infty(\Omega)\to L^\infty(\Omega;Y),\qquad b\mapsto bW,
\]
with \(\|M_W\|\leq L\).  Set
\[
  T=M_W\widetilde A.
\]
Then \(T:G\to L^\infty(\Omega;Y)\) is linear and \(\|T\|\leq L\).

Let \(X_n\) be simple \(S\)-valued random variables with \(X_n\to X\) in
\(L^p(\Omega;S)\) and almost surely.  Since \(\widetilde A\) agrees with \(A\)
on \(S\), the definition of \(T^{[p]}\) gives
\[
  (T^{[p]}X_n)(\omega)
  =J(X(\omega))(X_n(\omega))\,W(\omega)
  \quad\hbox{a.s.}
\]
Therefore
\[
  \|T^{[p]}X_n-Z\|_{L^p(\Omega;Y)}
  \leq L\|X_n-X\|_{L^p(\Omega;G)}\to0,
\]
because \(J(X(\omega))(X(\omega))W(\omega)=Z(\omega)\), with both sides zero
on \(\{X=0\}\).  Since \(T^{[p]}\) is bounded on \(L^p(\Omega;G)\), this proves
\(T^{[p]}X=Z\).
\end{proof}

\section*{Proof of Theorem \ref{thm:main}}

Apply Lemma \ref{lem:factor} with
\[
  G=\g(H,E_0),\qquad Y=\g(H,E_1),\qquad
  X=DF,\qquad Z=D\phi(F).
\]
The pointwise estimate proved above gives the hypothesis of the lemma.  The
resulting operator \(T\) is the required \(T_F\).

\section*{Novelty Check and Scope}

The argument answers exactly the operator-valued statement of
\cite[Proposition 3.10]{PV}: \(T_F\) lies in
\(\cL(\g(H,E_0),L^\infty(\Omega;\g(H,E_1)))\).  It deliberately does not claim
the stronger and different object \(T\in L^\infty(\Omega;\cL(\g(H,E_0),
\g(H,E_1)))\).  Remark 3.11 already notes that such random-operator-field
formulations have additional measurability subtleties.

The run indexes were searched for arXiv:1204.2946 and the relevant Malliavin,
UMD, Lipschitz-chain-rule, basis, and bounded-approximation-property terms.
Targeted web searches on 2026-06-23 found the source paper but no later exact
answer to Remark 3.11 or to the basis-free strengthening proved here.

\begin{thebibliography}{9}
\bibitem{PV}
M. Pronk and M. Veraar,
\emph{Tools for Malliavin calculus in UMD Banach spaces},
Potential Analysis 40 (2014), 307--344; arXiv:1204.2946.

\bibitem{Kelley}
J. L. Kelley,
\emph{Banach spaces with the extension property},
Transactions of the American Mathematical Society 72 (1952), 323--326.
\end{thebibliography}

\end{document}
